\documentclass[10pt]{article}
\usepackage[
  paperwidth=595.22pt,
  paperheight=842pt,
  margin=0pt
]{geometry}
\usepackage{fontspec}
\usepackage{xeCJK}
\usepackage{tikz}
\usepackage{xcolor}
\usepackage{array}
\usepackage{booktabs}
\usepackage{tabularx}
\pagestyle{empty}

\setmainfont{Arial}
\setsansfont{Arial}
\setCJKmainfont{SimSun}
\setCJKsansfont{SimSun}
\xeCJKsetup{CJKglue={\hskip 0pt}}

\definecolor{textblack}{RGB}{20,20,20}

\newcommand{\textboxat}[5]{%
  \begin{tikzpicture}[remember picture,overlay]
    \node[
      anchor=north west,
      inner sep=0pt,
      text width=#3,
      align=#4,
      text=textblack
    ] at ([xshift=#1,yshift=-#2]current page.north west) {#5};
  \end{tikzpicture}%
}
\newcommand{\term}[2]{#1 (#2)}
\newcommand{\hcell}[2]{\begin{tabular}[t]{@{}l@{}}\textbf{#1}\\[-1pt]{\fontsize{6.4}{7.0}\selectfont\textbf{(#2)}}\end{tabular}}
\newcommand{\hcellthree}[3]{\begin{tabular}[t]{@{}l@{}}\textbf{#1}\\[-1pt]{\fontsize{6.4}{7.0}\selectfont\textbf{(#2)}}\\[-1pt]{\fontsize{6.4}{7.0}\selectfont\textbf{(#3)}}\end{tabular}}
\newcommand{\rcell}[2]{\begin{tabular}[t]{@{}l@{}}#1\\[-1pt]{\fontsize{6.4}{7.0}\selectfont (#2)}\end{tabular}}

\begin{document}
\mbox{}

\textboxat{72pt}{38.4pt}{55pt}{left}{\fontsize{11.2}{13}\selectfont 1P5}
\textboxat{158pt}{39.0pt}{280pt}{center}{\fontsize{10.4}{12}\selectfont P5 Thermodynamics Lab: Background \& Prep}
\textboxat{408pt}{38.4pt}{134pt}{right}{\fontsize{11.2}{13}\selectfont Page 26 of 52}
\textboxat{72pt}{792.4pt}{280pt}{left}{\fontsize{10.2}{12}\selectfont P5 Thermo PrepWork 2526.docx}
\textboxat{402pt}{792.4pt}{140pt}{right}{\fontsize{9.4}{11.2}\selectfont Jan-2026}

\textboxat{72pt}{69pt}{470pt}{center}{%
{\fontsize{10.3}{12}\selectfont \mbox{Table 1 - \term{温度传感器对比}{Comparison of Temperature Transducers}}}%
}

\textboxat{72pt}{92pt}{470pt}{left}{%
{\fontsize{7.2}{9.0}\selectfont
\setlength{\tabcolsep}{1.5pt}
\renewcommand{\arraystretch}{1.20}
\begin{tabular}{@{}
>{\raggedright\arraybackslash}p{54pt}
>{\raggedright\arraybackslash}p{48pt}
>{\raggedright\arraybackslash}p{45pt}
>{\raggedright\arraybackslash}p{56pt}
>{\raggedright\arraybackslash}p{74pt}
>{\raggedright\arraybackslash}p{88pt}
>{\raggedright\arraybackslash}p{86pt}
@{}}
\toprule
\hcell{类型}{Type} &
\hcell{成本}{Cost} &
\hcell{输出}{Output} &
\hcell{线性度}{Linearity} &
\hcell{温度范围}{Temp. Range, $^\circ$C} &
\hcellthree{单个传感器量程}{Individual Sensor}{Range, K} &
\hcellthree{易实现精度}{Readily Achievable}{Accuracy, K} \\
\midrule
\rcell{热电偶}{TC} & \rcell{中等}{medium} & \rcell{低}{low} & \rcell{合理}{reasonable} & $-250$ to $2500$ & $\approx1000$ & $\pm1.0$ \\
\rcell{铂电阻}{PRT} & \rcell{高}{high} & \rcell{低}{low} & \rcell{好}{good} & $-200$ to $850$ & $\approx500$ & $\pm0.1$ \\
\rcell{热敏电阻}{Thermistor} & \rcell{低}{low} & \rcell{高}{high} & \rcell{差}{poor} & $-50$ to $250$ & $\approx100$ & $\pm1.0$ \\
\bottomrule
\end{tabular}
}%
}

\textboxat{72pt}{245pt}{470pt}{left}{%
{\fontsize{10.2}{15.0}\selectfont
{\fontsize{12.3}{15}\selectfont 3.3\quad \term{旋转速度测量装置简介}{Introduction to Rotational Speed Measurement Devices}}

\vspace{5pt}
\term{旋转速度}{rotational speed} 的测量虽然不如压力 (pressure) 或温度 (temperature) 测量那样常见，但仍然应用广泛。它适用于内燃机 (internal combustion engines)、涡轮机械 (turbomachinery，包括风力涡轮机和喷气发动机) 以及飞机姿态与运动 (aircraft attitude and motion) 等旋转机械。

\vspace{13pt}
对于只有一个旋转自由度 (1-degree-of-rotational-freedom) 的机器，速度通常用每分钟转数 (revolutions per minute, rpm) 表示，而不是用 Hz。相比之下，飞机的旋转运动具有三个自由度 (3-degrees-of-freedom, 3-DoF)，通常通过陀螺仪 (gyroscope) 以弧度每秒 (rad/s) 测量。为简洁起见，这里只讨论一自由度 (1-DoF) 旋转测量；多自由度陀螺仪测量留待后续课程学习。

\vspace{13pt}
如今，大多数速度测量系统使用\term{脉冲计数方法}{pulse counting methods}。在这类方法中，感应式或光学探头 (inductive or optical probe) 检测旋转特征经过探头的时刻，并产生相应的时间脉冲列 (train of pulses)。被检测特征取决于具体设计和应用，可以是齿轮齿、六角螺母平面，或轴上的蚀刻/阳极氧化图案。
}%
}

\end{document}
